% =====================================================================
%  R. Nielsen, ON DARKENING OF THE MIND, A CONTRIBUTION TO THE PSYCHOLOGY OF SUICIDES
%  In: For Idea and Reality.  A Journal, edited by R. Nielsen, Bj. Bjørnson and Rud. Schmidt.  Volume for 1869.
%  Copenhagen: Chr. Steen & Søn; printed by Louis Klein.  pp. 689-706.  (Om Sindsformørkelse, Bidrag til Selvmorderes Psykologi, in For Ide og Virkelighed, Aargang 1869.)
%  ENGLISH TRANSLATION, from transcription.tex (Danish).
%  GENERATED by ebook-method/for-idee-1869/tr_build.py --write from this template and the translators' parts in
%  ebook-method/for-idee-1869/tr/ (TRANSLATOR-BRIEF.md; terms: GLOSSARY.md, the titles and four key terms approved by
%  Hans 2026-10-13).  Conventions: TRANSLATION-PLAYBOOK.md.  Page numbers of the 1869 print in the margin.
% =====================================================================
\documentclass[12pt,a4paper,openany]{book}
\usepackage[utf8]{inputenc}
\usepackage[T1]{fontenc}
\usepackage{libertinus}
\usepackage{libertinust1math}
\usepackage[english]{babel}
\usepackage[protrusion=true,expansion=true]{microtype}
\usepackage{setspace}
\usepackage[margin=1.3in]{geometry}
\usepackage{fancyhdr}
\usepackage{hyperref}
\hypersetup{pdftitle={On Darkening of the Mind}, pdfauthor={Rasmus Nielsen}, hidelinks}
\setlength{\headheight}{15pt}
\pagestyle{fancy}
\fancyhf{}
\fancyhead[LE,RO]{\thepage}
\fancyhead[C]{\textit{R. Nielsen: On Darkening of the Mind (1869)}}
\setstretch{1.3}
\usepackage{marginnote}
\renewcommand{\thefootnote}{\arabic{footnote})}   % the print numbers each page's note ¹) (p. 553: *))
\newcommand{\opage}[1]{\marginnote{\footnotesize\textit{[#1]}}}
\newcommand{\essayhead}[2]{\begin{center}{\large #1}\ifx&#2&\else\\[0.3em]{\normalsize #2}\fi\\[0.5em]\rule{3em}{0.4pt}\end{center}\medskip\noindent}
\newcommand{\signature}[1]{\par\smallskip\hfill #1\hspace*{2em}\par}

\begin{document}

% the volume's title page (PDF 7)
\begin{titlepage}
\begin{center}
\vspace*{3em}
{\Huge For Idea and Reality.}\\[1em]
\rule{8em}{0.4pt}\\[1em]
{\large A Journal,}\\[0.8em]
{\small edited}\\[0.4em]
{\footnotesize by}\\[0.6em]
{\large R. Nielsen,}\\[0.4em]
{\large Bj. Bjørnson and Rud. Schmidt.}\\[1em]
\rule{8em}{0.4pt}\\[1em]
{\large Volume for 1869.}\\[1em]
\rule{8em}{0.4pt}\\[4em]
Copenhagen.\\[0.4em]
Chr. Steen \& Søn, Publishers.\\[1.5em]
{\small Printed by Louis Klein.}\\[3em]
{\small Translated from the Danish}
\end{center}
\end{titlepage}

\essayhead{On Darkening of the Mind,}{A Contribution to the Psychology of Suicides.}
% --- p. 689 ---
\setcounter{footnote}{0}%
\opage{689}If anyone were to raise, quite generally, a question such as this: who is really the greater criminal, the murderer pure and simple or the suicide, assuming that both have acted with full deliberation? then to sound sense the answer must seem not to be in doubt. There is in the murderer's mind such a brutal egoism, such a cruel contempt for the life of his neighbor, such a revolting defiance of divine and human laws, that society rightly sees in this bloody evildoer its sworn enemy. It is otherwise with the suicide. Such an unfortunate man did not, after all, become dangerous to society. He did not push himself forward by force; no, he withdrew in all silence. He did not seize the dagger in order to clear others out of the way and make room for himself; on the contrary, as one who was superfluous in life, he wished himself to go out of the way, to remove himself in all unobtrusiveness, so that there might be room for others. True, he presumptuously tore asunder his own thread of life, but who knows? --- a venomous sting of the serpent of death had perhaps already lodged deep in his heart. The notion of what a human being must suffer and struggle, before in the shipwreck of life he lets go of the last plank, the last hope, and chooses the last way out of despair, may awaken serious thoughts in the beholder.

% --- p. 690 ---
\setcounter{footnote}{0}%
\opage{690}In order to find out, in each case at issue, what may have been the motive in the individual suicide, one naturally attends to the closer circumstances, and there especially to the partly chance utterances, partly express declarations, which the deceased may have left behind. But it is precisely the remarkable thing that such statements not infrequently stand in so cutting a disproportion to what they were supposed to explain that they unmistakably point to something else, something hidden, which has stirred in the mind, perhaps unconsciously or at any rate only half-consciously. What this, the actual and true ground of the deed, may have been, is a dark riddle, a secret which the dead man took with him to the grave. On these terms the beholder then has to do, not directly with the departed, but with himself. The real dissolves for him into mere possibilities; the longer he stares down into the abyss of these possibilities, the more he exposes himself to dizziness; a thought rises up from the depths and makes the blood run cold, a dark thought with the question: could not the like also befall you? The more ideality there is in the nature of the beholder, the more natural it will be for him to spread a softening gleam even over the crudest suicide.

We praise what is human in such sympathy; but for that reason we dare not forget what a suicide is in itself, that it is a \emph{crime}. We dare the less forget it, as the propensity to this crime is, even to an alarming degree, widespread among us.

How frequent suicide is on average in Denmark, can be gathered from the ``survey of the proportion of
% --- p. 691 ---
\setcounter{footnote}{0}%
\opage{691}suicides to deaths in general in the age classes over 10 years, 1835 --- 1845,'' which Dr. med. Keyser (Dansk Folkeblad 1846) has communicated. The author adds: ``One sees from this that in all Denmark every sixtieth % print: Tresinstyvende (as printed)
of the dead over ten years of age is a suicide, or in other words that almost two percent of all deaths are caused by this; that this frequency rises in Copenhagen; that for males in this city it reaches an alarming height, in that every twenty-seventh male dying over ten years of age is a suicide, and that thus in this age almost four percent of the deceased are carried off in this way.'' This fact must greatly astonish anyone who has not himself had occasion to investigate the matter; it must, as the author indeed also emphasizes, recall the great difficulty that is connected with making a reliable experience, and ``how much mischief is done with this word.''

In ``\emph{Stages on Life's Way},'' a work that appeared in 1845, the pseudonym expresses his gladness at being an author in a small country (page 378), and says: ``Some of my countrymen no doubt think that Copenhagen is a tedious city and a small city. To me, on the contrary, refreshed by the sea by which it lies, and without, even in winter, being able to give up the memory of the beech wood, it seems to be the happiest place of abode that I know. Large enough to be a larger city, small enough that there is no market price on human beings; can the tabular consolation, such as one has it in Paris, over so and so many suicides\dots{}. not intrude disturbingly'' etc. Although we now willingly concede that poetic moods and statistical tables are
% --- p. 692 ---
\setcounter{footnote}{0}%
\opage{692}incommensurable magnitudes, we nevertheless remain convinced that the allusion to the tabular consolation over so and so many suicides would have turned out differently if the author of the ``Stages'' had had any notion of the result to which a ``comparative survey of the frequency of \emph{suicide} in Denmark, 1835---44, and in France, 1835---43,'' has actually led; for the result is that \emph{suicide is about three times as frequent in Denmark as in France,} and that the \emph{frequency} of suicide, both for each sex singly and for both together, is \emph{greater} in \emph{Copenhagen} than in the whole \emph{department of the Seine}.

He who inquires into the universal psychological basic condition of suicide meets with great difficulties everywhere. To be sure, difference in age and sex may exert an influence, but what influence? The starting point that is sought in the differences of temperament is highly unreliable, for experience teaches that the sanguine and the phlegmatic just as well as the choleric and the melancholic can be found on the list of suicides. No difference between the sick and the healthy, between the cowardly and the brave, between the cultivated and the uncultivated, between the rich and the poor, indeed not even between the exceedingly happy and the boundlessly unhappy can be established with due certainty on this enigmatic point.

That a man of defiant courage throws life away when it becomes a burden to him, one finds natural, but that cowardice and anxiety can lead to the same result must seem strange. Yet there are examples. Just as for many it is a mortally dangerous matter to sit at table self-thirteenth, so too it is regarded by the common folk
% --- p. 693 ---
\setcounter{footnote}{0}%
\opage{693}in the country as mortally dangerous for one to fall down into the grave during the throwing of earth; for he who falls down to the dead must die before the year is out. A man whom the mishap befell took it so much to heart that he hanged himself a couple of months later: from fear of death he chose to die.

A suicide may be committed in rage; but that it can also be carried out in perfect composure, with a clear view of the situation, is attested by the following letter, with which a peaceable lodger wished to calm his landlady: ``Dear Madam N.! In the state in which you will find me, you will at once cry out and scream: he has hanged himself, but that you must not do. You need only say, I am dead\dots{} I am having myself buried from the Hospital. Mr. H. is to have some money, tell him, it will be duly paid by one of my friends, from whom I have borrowed 40 rix-dollars. (Folkebl. 1846, p. 72). This phlegmatic imperturbability stands in cutting contrast to the bitterness with which the craftsman's apprentice (loc. cit., p. 75), who had received a box on the ear, took leave of the world and wrote to his parents: ``For the last time I will say farewell to you here in life; I cannot bear it any longer. Curse Madsen, curse the journeyman; two beatings in one night.''

How cultivation and coarseness can, despite the difference of motives, meet in the same crime, is a well-known matter. In the second part of \emph{Either --- Or} (second edition, p. 172) the ethicist sketches a psychological picture of life, apt from his standpoint, of a young man whom intellectual overstrainedness and brooding inwardness drove to
% --- p. 694 ---
\setcounter{footnote}{0}%
\opage{694}suicide. ``Should anyone,'' writes the deceased in his letter to the Councillor of Justice, ``ask you about the occasion, you may say that it was on the occasion of the great fire. Should anyone ask you about the time, you may say that it was in the month of July, so remarkable to me\dots{} I do not regard a suicide as anything praiseworthy. It is not from vanity that I have resolved upon it. On the contrary, I believe in the correctness of the proposition that no human being can bear to see the Infinite. It has shown itself to me once in the intellectual respect, and the expression for it is ignorance. Ignorance is, namely, the negative expression for infinite knowledge. A suicide is the negative expression for infinite freedom. It is a form of infinite freedom, but the negative form. Happy he who finds the positive one\footnote[1]{Perhaps the story is invented. But that the difference between fiction and reality means nothing where the talk is only of \emph{psychological truth}, is a remark that he \emph{who} writes for cultivated readers may well regard as superfluous.}.'' To % print: note mark „¹;«.“ (as printed)
this striving for freedom in a dialectically refined form of cultivation one finds (Folkebl. p. 74) a counterpart in the form of undialectical coarseness. It is a locksmith's apprentice who longed for freedom in America, but, for want of a passport, had to seek it in one of the canals of Frederiksberg Garden, on whose bank he left the following: ``Here lies the undersigned right opposite. Since I have left my master, because I would not crawl before the mistress, and now when I intended to travel to America, he would not write on my certificate, as he ought to, I could therefore get no passport. What was I to do now. I knew only this counsel \dots{}''

% --- p. 695 ---
\setcounter{footnote}{0}%
\opage{695}To what degree exalted deification of genius, fantastic cult of the genius, platonic, hypertrophic infatuation with the idea can attack female nerves, Charlotte Stieglitz, otherwise admirable for her spirit and grace, affords a tragic example\footnote[1]{Cf. ``A Drama in Literary Life'' (by Taillandier. Tales and Sketches from the Feuilleton of ``Fædrelandet.'' Translated by Rudolf Schmidt. P. 339 ff.}. She regarded her husband as % print: (af Taillandier. (as printed)
the poet-Messiah, from whose brilliant appearance all Germany --- and thus the whole world --- was to await its poetic redemption; it was therefore not enough for her to be his wife, she had also to be his Muse. The expectations were disappointed, the genius failed, the man pined away; the wife chosen as Muse now placed her only hope in this, that a strong shock to the mind might perhaps give back to the sick man his lost elasticity and awaken his slumbering genius; in order to bring about such a spiritual shock and to consecrate her poet by the baptism of suffering, she adorned herself as the bride of Death and then thrust the dagger into her heart.

The shaking of the soul by which Charlotte Stieglitz wished to lift her idea-sick spouse up to the glory of the ideal stands in cutting contrast to the crushing which a vindictive servant girl by her sacrifice has evidently intended for her faithless lover. ``Mr. Sivertsen! Now you can see what you have got for your wicked actions against me, who 15 weeks ago suffered so much for your sake. And you will show yourself so brazen as to say that you are not the father of my child, as God in my hour of death shall be witness; but you are a wicked person, with whom
% --- p. 696 ---
\setcounter{footnote}{0}%
\opage{696}it can never go well, for you do so much evil to your fellow human beings. You gave me fine promises, but I did not see much of them. Farewell, wicked person.'' (Folkebl. p. 75).

At a cursory glance it might almost seem unfitting to place side by side two beings so unlike as that enthusiast intoxicated with ideas and this servant girl, by no means rich in spirit; and yet, when the matter is to come before the court of psychology and ethics, it is a great question whether the aberration of the idea does not, in the sense of eternity, have far more dangerous consequences than the raw straying. Where the motives are so different, the actions too must surely be different. But the dissimilarity of the actions is lost here in the similarity of the deed: it is in reality the same deed, the same crime.

The commonest motives to suicide Dr. Keyser (Folkebl. p. 52) has brought under the following seven categories: cares of livelihood, unhappy love, family cares, bodily illness, dissatisfaction with one's position and with life, bad conduct and weakness of mind. Insofar as this classification follows factual, in particular French, data, it may well have its validity in tabular respects; but from a psychological standpoint the category: dissatisfaction with life, is obviously too general, and the category: weakness of mind, when the word is taken in its usual meaning, so special that it excludes accountability. The motives cited become remarkable chiefly by the way in which they combine with the influence of the seasons.

That the change of the seasons, like the alternation of day and night, affects the mood of the mind, is a
% --- p. 697 ---
\setcounter{footnote}{0}%
\opage{697}familiar matter. One would now think that the light and warm part of the year, richest in external conditions of life, must count the fewest. the dark, cold season, weighed down by cares of livelihood, the most suicides; but, strangely enough, the relation is just the reverse: % print: færreste. den (as printed)
December and January have the fewest, June and July the most. Just as enigmatic is the influence of the seasons on the manner of death. We can at a pinch understand that hanging is on the whole preferred in winter, since the water is so cold, and that drowning has something alluring in the heat of summer; but what is inexplicable lies in the transition. ``After hanging in the month of December has sunk to its lowest, it rises suddenly to its highest in January, and after drowning has reached its highest in December, it sinks suddenly to its lowest in January.'' If this is more than a play of chance, then the causes must be so hidden that no one can yet divine in what they consist; only so much can be said, that no rule forms itself without a circle of standing conditions, and that no effect is without cause.

If we seek a unitary expression for the condition of soul in which the suicide must be assumed to find himself at the moment of the deed, then there is hardly a word either more comprehensive or more significant than \emph{darkening of the mind.} Just as the darkness of night makes the eye blind and lets everything run into one, so too the darkness in the soul can let the most various motives run together; the light of reason can then in a way be said to have gone out before the physical life is extinguished. The point on which the psychological interest fastens
% --- p. 698 ---
\setcounter{footnote}{0}%
\opage{698}is the origin of the darkening of the mind, its slow increase, its sudden onset.

Whence comes darkening of the mind? If we disregard the weakness that must chiefly be ascribed to organic disturbances in the brain, then it may well be assumed that it comes from within, from the ground of the mind, from a secret darkness that has withdrawn itself from consciousness. Just as the whole globe is never lit at one time, but one half always has night while the other has day, so too the life of the soul is divided between light and darkness. Yet the comparison limps; for while the globe, by turning on its axis, is successively lit on all sides, there is one side of the life of the soul, the organic, namely, where it is always night; and while the interior of the earth excludes the sunlight for all time, the interior of the soul is precisely destined to open itself to the light of consciousness. The interior of the soul is the self, and self-consciousness consists precisely in this, that a human being makes his original self clear to himself. The darkness that forms the natural boundary of the life of consciousness belongs to our organization and strikes no fear; but the darkness in the self, the spiritual darkness, which withdraws from consciousness and yet is destined to yield to consciousness, is like an omen of the underworld, that night of life which harbors the dangerous possibilities.

The spiritual darkness can brood in the mind, and yet it may seem as if it were broad daylight in consciousness. The world of things with its phenomena, existence with its tasks and finite interests, the whole external reality can stand before intuition in full clarity, the understanding can be highly developed and well oriented in many directions, and yet is the darkness
% --- p. 699 ---
\setcounter{footnote}{0}%
\opage{699}within. Thus such a human being can live with others, work and enjoy like others, fear and hope, laugh with the glad and weep with the sorrowful: there is no one in his surroundings who notices anything; indeed the person concerned himself has hardly any notion of how his condition is, that spiritually speaking he goes with death in his heart. Now and then he has perhaps sensed an involuntary prompting; it was a vision of terror, a will-o'-the-wisp of the imagination, a wild thought that came and vanished; but that was all. This is the disposition.

Of this disposition holds true what S. Kierkegaard in his work \emph{The Sickness unto Death}, p. 17, states about despair: ``As the physician must no doubt say that perhaps there lives not a single human being who is entirely healthy, so one would have to say, if one really knew the human being, that there lives not a single human being who is not nevertheless a little in despair, in whom there does not nevertheless dwell, deep within, an unrest, an unpeace, a disharmony, an anxiety about an unknown something, or about something he does not even dare to make acquaintance with,\dots{} an anxiety about himself, so that he nevertheless, as the physician speaks of going about with a sickness in the body, goes about with a sickness, goes about and carries a sickness of the spirit, which once in a single while, in glimpses, by and with an anxiety inexplicable to himself, lets itself be noticed, by the fact that it is in there.'' The disposition is common, so common that it can even be assumed to be natural to the natural human being; but to what degree it is present in individuals, whether the darkness of the soul is to vanish before the light of the spirit, or whether, as a secret of the night, it is some day in an unhappy moment to break forth from its hiding place and darken the whole mind, will depend on how a human
% --- p. 700 ---
\setcounter{footnote}{0}%
\opage{700}being deals with himself, and besides rest as much on external circumstances as on inner conditions.

To drive the spiritual darkness out of the soul two conditions are required: essential self-understanding and a character formation firm of will. Where one of these conditions is lacking, the disposition may easily, under given circumstances, become so strong that it ends in suicide.

That the conditions under which alone the natural and yet so unnatural disposition can be overcome are only too often lacking, has its innermost and essential ground in self-neglect, and this in turn in self-forgetfulness. Self-forgetfulness has its degrees; it can be partly immediate and naive, partly moral, partly religious in the strictest sense.

Naive self-forgetfulness is characterized by this, that the individual simply apprehends his life as a gift of nature without considering that this gift of nature is at the same time a spiritual task. He who is truly to learn what life is must not merely have something to \emph{live} on, he must in the unconditional sense also have something to \emph{live} for; the latter is just as necessary as the former. To the natural sense, however, the matter presents itself differently. ``If only we had enough to \emph{live} on, we could easily get something to \emph{live} for; for why do we \emph{live}? We are to work, to act, and to do good! that is, to be sure, life's law, its strict condition; but it is not life's aim. The aim is to rejoice in life, to enjoy life. Only he knows how to live who understands how to enjoy: the greater the enjoyments, the more the pleasures, the more beautiful, richer and happier is life.'' That this is a
% --- p. 701 ---
\setcounter{footnote}{0}%
\opage{701}very widespread and popular view, no one will surely deny\footnote[1]{How natural the point of view. that life's goal is to enjoy and amuse oneself. can take root as a dogma, and how involuntarily this dogma can come to expression even among the elderly, of this we shall merely cite a little example. A couple of respectable burghers sat together one evening in the theater and saw Don Juan. The play is over; the lightning bolt has struck. The curtain falls. Then the one says: ``It was a cruel end, though, that he came to.'' ``Yes,'' answered the other, ``but he had amused himself, after all.''}. % print: Synsmaade. at, sig. kan, ramt. Teppet (as printed)

To be sure, one gets now and then, as it were, an inkling that life might possibly have another and higher aim; but on the first natural stage of self-forgetfulness that does not signify much. Indeed, if it were a matter of phrases, one could easily gain a number of high-flown reflections; but life is a reality, and in reality what then is life? ``A breath in the rushes,'' a dark riddle. Since nothing comes of sitting down to brood over such a riddle, one takes the world as it is, enjoys as long as one can, and amuses oneself.

When then life weighs heavily, when the prospects darken, when he who relied on fortune, who strove after enjoyments, is struck by misfortune and oppressed by adversity, what wonder, then, that ``cares of livelihood,'' ``unhappy love,'' ``family cares'' and the like can in a fateful moment overwhelm the unprepared, chase the darkness up from the mysterious night and drive him to despair! The more the view that life exists for the sake of enjoyment spreads and takes root among the common people, the less may we wonder that suicide is so frequent, we must rather wonder that it is not yet more frequent.

% --- p. 702 ---
\setcounter{footnote}{0}%
\opage{702}The more the frenzy for enjoyment spreads and sets the masses aflame, the more spiritual immaturity grows. At what stage of immaturity naive self-forgetfulness can hold individuals captive is evident when we consider the influence that the seasons exert on the act that decides between life and death. It is with suicides as with birds of passage: it is not a will matured to freedom, but an obscure natural power, that moves them to depart. That it happens in the reverse order, that the birds of passage prefer to stay here in the summer and only depart when winter approaches, while suicides as a rule prefer to winter over in this life and only in greater numbers depart from here in the summer, gives a faint indication that the human being is spirit. It is, namely, a psychological fact that summer is more dangerous for the immature children of the spirit than winter. The hard necessity, the momentary want, spurs the instinct of self-preservation, calls to exertion and drives away whims; whereas the milder season, smiling nature with its soothing, bewitching impressions, tempts the idle in particular to torpor, daydreams and confusing thoughts.

To interrupt self-forgetfulness and awaken the thoughtless, existence does what it can; as rain alternates with sunshine, calm with storm, so it lets fortune change to misfortune, prosperity to adversity, joy to sorrow; it comes to meet the human being, speaks through the dispensations, and says to the confused: go into yourself! Those who are now awakened by such admonitions gain thus far a higher notion of the seriousness of life, as they recognize that it is will and character on which
% --- p. 703 ---
\setcounter{footnote}{0}%
\opage{703}it depends. Life is a school for character. At this stage much capability and endurance can unfold, much civil and political activity. Here is something to live for; here are rational aims. A vigorous self begins to form. It is, as one sees, a first approach to the ethical; character stands, so to speak, in the forecourt of ethics, but yet, despite its apparent strength and bearing, on weak feet, if the mind is one-sidedly turned outward, exclusively occupied with life in the external. The whole question of the self in the deeper sense then becomes, as S. Kierkegaard puts it, ``a kind of blind door'' in the background of the human soul, inside which there is nothing. The individual ``takes over what he in his language calls his self, that is to say, whatever abilities, whatever talents, etc., may have been given him, all this he takes over, yet with the direction outward, toward, as it is called, life, the real, the active life; he handles very cautiously the little bit of reflection he has in him, he fears that it will come up again, That in the background. Then he succeeds little by little in forgetting it; in the course of the years he finds it almost laughable, especially when he is in good company with other capable and active men, who have a sense and capability for real life. Charming! He is now, as it says in the novel, already for several years happily married, an active and enterprising man, father and citizen, perhaps even a great man; at home in his house the servants call him ``He Himself''\footnote[1]{The Sickness. unto Death. P. 53---54.} --- and yet there is in this % print: Sygdom. til Døden. (as printed)
``He Himself,'' when ethics conducts its inspection,
% --- p. 704 ---
\setcounter{footnote}{0}%
\opage{704}a bottomless obscurity, a selfless darkness: the eternal light is lacking. On the dark, selfless ground a selfish, firm will can nevertheless raise a structure of character that in certain ways resembles the Egyptian temples: the approach is grand, the masonry imposing, the colonnaded halls splendid, but the Holy of Holies is a little dark hole. If such an individuality now carries the unhappy disposition in itself, then it takes only a blow of fate, an offense he cannot tolerate, a loss of money he cannot bear, a hasty step that cannot be undone, a bodily weakness he will not put up with, and what has he now to live for? ``The strong self is its own master; a rope about the neck, a bullet in the forehead, a leap into the sea, what is that, then?'' % print: mere?» (as printed)
So it happens. The surroundings, which suspect nothing, suddenly hear an unexpected message, and ask in surprise: who would have thought that such a life should have such an end?

The ethical self-neglect has its deeper ground in the religious. We have in the preceding treatises sought to point out the principled connection between religion and ethics. It is a bewitching conceit that the human will, without a living impression of God's will, should be able to do justice to the idea.

The more ideality there is in a human being, the more easily he is exposed to aberrations of the spirit when religious self-understanding is lacking. It is in this respect quite significant that the heroine of that sinister ``Drama'' had already strayed away from the religious before she went astray in ideas. As a quite young girl Charlotte Willhoeft, under the influence of a teacher who
% --- p. 705 ---
\setcounter{footnote}{0}%
\opage{705}belonged to the sect of the Methodists, had given herself over to a fanatical and dark piety; she despised this world, she mocked its active life, and more than once, in her striving toward God, the thought of suicide passed through her brain. This glowing longing for death, this impatient striving toward another life, was, however, only a crisis; she returned to earth and ``preserved from her communion with God only a love for everything that is divine among human beings.'' This is then an eloquent proof that religious fancies of mood, without self-knowledge, without deeper understanding of the life of the will and its tasks, are so far from clearing up the essence of the self that the wildly crossing flashes of light, on the contrary, intensify the spiritual darkness.

It is especially loose notions of God's grace, of the forgiveness of sins and eternal blessedness that confuse the weak. ``I am now journeying,'' writes a suicide (Folkebl. p. 74), ``yonder to eternal rest, where I, if I find God's grace, after having prayed for forgiveness or pardon, shall come to rejoice and be glad with a higher all-governing Being there, where no sin and no wrong takes place\dots{} I shall then look down from on high on you all here on earth''\dots{} There is no doubt that this somewhat soft, but, as it seems, at the same time very childlike suicide really did have religious moods; he has surely also, when it came to the point, ``known his Scripture passages.'' But it is one thing to know Scripture passages, another to practice religious thoughts of life. As regards Scripture passages, it does not depend on the quantity, but on the appropriation. A single passage,
% --- p. 706 ---
\setcounter{footnote}{0}%
\opage{706}--- it is found in Paul's Epistle to the Romans, chap. 14, vv. 7-8 --- when it is grasped in living faith and appropriated in spirit and truth, is a reliable shield and sure bulwark against the temptation to suicide.

If a lively occupation with biblical notions and ``godly meditations'' were a reliable mark of true religiosity, it would still not be so difficult to go into battle against the darkening of the mind and to get the number of suicides somewhat reduced; but the religious task of life is in the strictest sense a task of the self\footnote[1]{Those who, on the occasion of the latest humane movements in the metaphysics of abstraction, possibly assume that there is really something afoot here, that the fullness of time has come when God must renounce having a will, and that man, who for his part naturally voluntarily gives up the immortality of the soul, has only, in infinite resignation, to reconcile himself with the Abstract, would, before they too credulously immerse themselves in the bouillon drawn from Spinoza, Strauss and Feuerbach, thinned by repeated reheating and to that extent cheap, --- for caution's sake --- do well to take up the question of the eternity of the self once more, and, for example, to examine whether criticism, which otherwise so pressingly refutes everything, has really ``refuted'' what S. Kierkegaard in his work ``The Sickness unto Death'' has elucidated about the sick self.}. Those who do not know what the self is
do not know either what faith is. To become aware of the self, its obscure intimations, its secret torments, its crises of sickness and its healing, is the task. But this is also the most difficult of all tasks; for there is nothing in the world that the natural human being, precisely because of the selfish, more gladly flees than just his own self.

\signature{R. \emph{Nielsen}.}

\end{document}
